
\textbf{Corollaire 6.8. (Grothendieck).} \emph{Soient @@M0@@ projectif sur @@M1@@ noeth\'erien, @@M2@@ tr\`es ample pour @@M3@@, et @@M4@@ une famille de classes de faisceaux coh\'erents sur les fibres de @@M5@@. On suppose}~:
\begin{enumerate}[label=(\alph*)]
\item \emph{Il existe un faisceau coh\'erent @@M6@@ sur @@M7@@ tel que @@M8@@ soit contenue dans la famille des classes de quotients de faisceaux de la forme @@M9@@.}
\item[@@M10@@] \emph{Dans les polyn\^omes de Hilbert @@M11@@ des @@M12@@, les coefficients en degr\'es @@M13@@ restent major\'es.}
\end{enumerate}
\emph{Alors les @@M14@@ (lorsque @@M15@@ reste dans @@M16@@) forment une famille limit\'ee ; de plus les coefficients dans @@M17@@ en degr\'es @@M18@@ restent minor\'es.}

En effet, on peut \'evidemment supposer @@M19@@ et @@M20@@. Alors la premi\`ere assertion r\'esulte formellement de 6.7 (ii) et 1.13 ; la deuxi\`eme r\'esulte en raisonnant par r\'ecurrence \`a l'aide d'une suite exacte @@M21@@, de 6.4 (ii).

\textbf{D\'efinition 6.9.} \emph{Soit @@M22@@ un corps. On appelle \emph{@@M23@@-cycle positif sp\'ecial de dimension @@M24@@} tout @@M25@@-sch\'ema projectif @@M26@@, muni d'un faisceau @@M27@@ tr\`es ample, qui est r\'eunion de sous-sch\'emas ferm\'es @@M28@@ de dimension @@M29@@, @@M30@@ \'etant obtenu par un changement de base @@M31@@ comme image inverse d'un @@M32@@-sch\'ema \emph{int\`egre} @@M33@@ (muni de @@M34@@). On appelle \emph{degr\'e} de @@M35@@ le coefficient @@M36@@ dans le polyn\^ome de Hilbert @@M37@@.}

\textbf{Lemme 6.10.} \emph{Soient @@M38@@ un corps alg\'ebriquement clos et @@M39@@ un @@M40@@-cycle positif sp\'ecial de dimension @@M41@@ et degr\'e @@M42@@. Alors @@M43@@ est un @@M44@@-faisceau.}

